\documentclass{beamer}
% \documentclass[handout,xcolor=pdftex,dvipsnames,table]{beamer}

\mode<presentation>
{
  \usetheme{Warsaw}
  \usecolortheme{whale}
  % or ...

  \setbeamercovered{transparent}
  % or whatever (possibly just delete it)
  \setbeamertemplate{navigation symbols}{}
}

\setbeamertemplate{itemize items}[ball]
\setbeamertemplate{itemize subitem}[triangle]
\setbeamertemplate{itemize subsubitem}[circle]
\setbeamercovered{transparent}

\usepackage{palatino} 
\usepackage{listings} % Gives syntax highlighting for python code. 
\usepackage{color} % Used for syntax highlighting. 
\usepackage{textcomp} % Used for syntax highlighting. 

% This gives syntax highlighting in the python environment 
\definecolor{gray}{gray}{0.5} 
\definecolor{key}{rgb}{0,0.5,0} 
\lstset{
language=python,
basicstyle=\ttfamily\tiny, 
otherkeywords={1, 2, 3, 4, 5, 6, 7, 8 ,9 , 0, -, =, +, [, ], (, ), \{, \}, :, *, !}, 
keywordstyle=\color{blue}, 
stringstyle=\color{red},
showstringspaces=false,
alsoletter={1234567890},
otherkeywords={\ , \}, \{},
keywordstyle=\color{blue},
emph={access,and,break,class,continue,def,del,elif ,else,%
except,exec,finally,for,from,global,if,import,in,is,%
lambda,not,or,pass,print,raise,return,try,while},
emphstyle=\color{black}\bfseries,
emph={[2]True, False, None, self},
emphstyle=[2]\color{green},
emph={[3]from, import, as},
emphstyle=[3]\color{blue},
upquote=true,
morecomment=[s]{"""}{"""},
commentstyle=\color{gray}\slshape,
emph={[4]1, 2, 3, 4, 5, 6, 7, 8, 9, 0},
emphstyle=[4]\color{blue},
literate=*{:}{{\textcolor{blue}:}}{1}%
{=}{{\textcolor{blue}=}}{1}%
{-}{{\textcolor{blue}-}}{1}%
{+}{{\textcolor{blue}+}}{1}%
{*}{{\textcolor{blue}*}}{1}%
{!}{{\textcolor{blue}!}}{1}%
{(}{{\textcolor{blue}(}}{1}%
{)}{{\textcolor{blue})}}{1}%
{[}{{\textcolor{blue}[}}{1}%
{]}{{\textcolor{blue}]}}{1}%
{<}{{\textcolor{blue}<}}{1}%
{>}{{\textcolor{blue}>}}{1},%
numbers=none,
}

\title[]{Python Workshop\\
Numpy Arrays}

\author[Langevin] % (optional, use only with lots of authors)
{C.D.~Langevin}
\institute[USGS] % (optional, but mostly needed)
{
  U.S. Geological Survey\\
  Reston, Virginia, USA
  }

\titlegraphic{\includegraphics[scale=0.5]{figures/c_USGSid1.pdf}}

\date[UQ12] % (optional, should be abbreviation of conference name)
{USGS National Groundwater Workshop, August 2012}

\subject{Python}


\begin{document}

\begin{frame}
  \titlepage
\end{frame}

%put this pdf image on all following slides
\logo{\vspace{-0.3cm} \includegraphics[width=1.5cm]{figures/c_USGSid1.pdf}\hspace*{11.10cm}}  

\begin{frame}{Outline}
\tableofcontents
\end{frame}

\section{Numpy Arrays}
\begin{frame}[fragile]
\frametitle{What is Numpy}
\end{frame}

\begin{frame}[fragile]
\frametitle{Creating an Array}
\end{frame}

\begin{frame}[fragile]
\frametitle{Basic Operations}
\end{frame}

\begin{frame}[fragile]
\frametitle{Arrays and Functions}
\begin{lstlisting}
In [55]: T = 2500 #ft2/d

In [56]: S = 0.01 #unitless

In [57]: t = 1.0 #d

In [58]: Q = 10000 #ft3/d

In [59]: r = concatenate( (array([1]), arange(10,110, 10)) )

In [60]: r
Out[60]: array([  1,  10,  20,  30,  40,  50,  60,  70,  80,  90, 100])

In [61]: u = r ** 2 * S / 4. / T / t

In [62]: from scipy.special import expn

In [63]: wu = expn(1, u)

In [64]: s = Q / 4. / pi / T * wu

In [65]: s
Out[65]:
array([ 4.21388046,  2.74804077,  2.30686505,  2.04889705,  1.86597563,
        1.72420422,  1.60848432,  1.51076188,  1.42622927,  1.35178543,
        1.2853129 ])
\end{lstlisting}
\end{frame}


\begin{frame}[fragile]
\frametitle{Universal Functions}
\end{frame}

\begin{frame}[fragile]
\frametitle{Indexing, Slicing, and Iterating}
\end{frame}

\begin{frame}[fragile]
\frametitle{Shape Manipulation}
\end{frame}

\end{document}